This note reports a measurement of dimuon production in \PbPb\
  collisions at \snn=5.02~\TeV\ through the process \gamgammumu,
  stimulated by the interaction of the large electromagnetic
  fields of the incident nuclei.
The measurements are performed using data from teh 2015 and 2018
  heavy-ion runs corresponding to an integrated luminosity of
  1.98~\inb.
This process has been considered previously in ultra-peripheral
  collisions, where the nuclei do not interact hadronically;
  the measurement reported here considers
  this photon-induced dimuon production in a broader class of \PbPb\
  collisions including those with hadronic interactions.
Sources of background such as heavy flavor decays were removed from
  the analysis through a template fitting method, which using
  \dZeroPair, the combined $d_0$ of the two muons as a discriminant.

The muons are observed have nearly equal and opposite transverse
  momentum resulting in steeply falling Acoplanarity
  (\Aco) and \kperp\ distributions.
However, the shape of the dimuon Acoplanarity distribution is
  observed to change significantly from
  UPC$\longrightarrow$peripheral$\longrightarrow$
  to central collisions.
It particular, for the distributions in central collisions,
  the peak of the distribution is shifted to non-zero values
  of \Aco, and \kperp.
At fixed centrality, the change in shape for the \Aco\
  distribution is seen to have a strong \pt\ dependence,
  but the shape of the \kperp\ distribution is quite stable
  with \pt.
The shift of the peak in the \kperp\ distributions are
  quantified using multiple functional forms.
The quantification of the shift gives a measure of the momentum
  kick(s) that the muons get from interactions with the QGP,
  or with the EM-fields of the colliding ions.
Apart from measuring the \Aco\ and \kperp\ distributions,
  the \NormYields; which measure the fraction of the total
  yield of such \gamgammumu\ pairs, per centrality percentile
  are also measured.
The \NormYields\ are observed to increase from peripheral to
  central events with the increase being larger at higher \pt.


Dimuon production in UPCs is under good theoretical control.
Models based on the equivalent photon approximation and leading-order QED
  are accurate in describing the dimuon invariant mass, rapidity
  and angular distributions.
As the physical picture underlying these models does not accommodate a
  significant impact parameter dependence to the acoplanarity distributions,
  the results presented in this note suggest an additional physical mechanism is
  responsible for the observation of the centrality-dependent broadening
  and change in the peak location of the acoplanarity distributions.
One possible mechanism, consistent with established phenomenology of the matter
  produced in heavy-ion collisions, is the QED analog of jet quenching.
If this mechanism is effective, dimuons from photon-photon reactions represent
  a new tool to study the microscopic structure of the quark gluon plasma.
